\section{Coverage and Control Implications}
\subsection{Coverage under score-distribution shift}
Let $P$ be the calibration score law and $Q$ the test score law. The model
and score function are fixed before collecting calibration data. We use
$\TV(P,Q)=\sup_A|P(A)-Q(A)|$.

\begin{theorem}[Shifted one-step coverage]
Let $S_1,\ldots,S_n$ be iid from $P$ and let $S_*$ be an independent
test score from $Q$. If $\TV(P,Q)\leq\epsilon$, then the threshold
in~\eqref{eq:quantile} satisfies
\begin{equation}
 \Pr\{S_*\leq q_\alpha\}\geq 1-\alpha-\epsilon.
 \label{eq:tv}
\end{equation}
The probability averages over both calibration data and the test score.
\end{theorem}
\begin{proof}
Introduce an independent $S'\sim P$. Exchangeability gives
$\Pr(S'\leq S_{(k)})\geq k/(n+1)$ for $k\leq n$, with ties
making the inequality conservative. For every realized threshold $q$,
$Q(({-\infty},q])\geq P(({-\infty},q])-\epsilon$.
Taking expectation over calibration proves~\eqref{eq:tv}. If $k>n$,
the infinite threshold covers every finite score.
\end{proof}

The shift parameter has an RL interpretation. Let $\mu_P,\mu_Q$ be
state--action laws and $K_P,K_Q$ their conditional next-state kernels.
Contraction of TV through the fixed score map and the triangle inequality
give
\begin{equation}
 \begin{split}
 \TV(P_S,Q_S)\leq{}&\TV(\mu_P,\mu_Q)\\
 &+\sup_{s,a}\TV(K_P(\cdot\mid s,a),K_Q(\cdot\mid s,a)).
 \end{split}
 \label{eq:decomposition}
\end{equation}
The first term captures occupancy changes induced by the policy; the second
captures changes in physical dynamics. The bound motivates evaluating both
policy families and physical conditions. In deterministic simulators its
kernel term can be one, so residual measurements provide a more informative
empirical diagnostic than treating~\eqref{eq:decomposition} as a numerical
certificate. Robust predictive inference develops stronger distributional
assumptions for such transfers \cite{cauchois2024robust,barber2023beyond}.

\subsection{From residual envelopes to safety margins}
If complete calibration episodes are iid, their maxima $M_H$ are iid even
when transitions within an episode are dependent. Applying Theorem~1 to
the maximum-score laws gives
\begin{equation}
 \Pr\{M_H(\tau_*)\leq q_H\}\geq1-\alpha-\epsilon_H,
 \label{eq:trajectory}
\end{equation}
where $\epsilon_H$ bounds TV between calibration and test maximum-score
laws. This construction accounts for temporal dependence at the sampling
unit at which the envelope is calibrated.

For a feasible action and any retained hypothesis $\theta$, Lipschitz
continuity yields the deterministic implication
\begin{equation}
 h(s_{t+1})\geq h(f_\theta(s_t,a_t))
       -L_h\|s_{t+1}-f_\theta(s_t,a_t)\|_2.
 \label{eq:lipschitz}
\end{equation}
Thus a residual of at most $m_t/L_h$ implies nonnegative next-state margin.
A probabilistic safety result requires coverage of the applied-action and
retained-hypothesis residual law, together with feasibility. Independent-score
coverage alone does not establish these closed-loop conditions.
Our evaluated coefficient $\eta=1$ is an empirical operating point,
distinct from the analytical $L_h=\sqrt{3.25}$; the experiments measure its
realized safety and control cost.
